\chapter{От семантического gate к write-path: \texttt{exp4} и \texttt{exp5}}

\section{\texttt{exp4}: semantic gate как первая попытка исправить retrieval}

\subsection{Что именно проверялось}

После \texttt{exp3} гипотеза выглядела почти естественно: если диффузия никак не
учитывает релевантность узла запросу, то, может быть, достаточно встроить в нее
множитель

\[
\mathrm{rel}(B\mid q),
\]

зависящий от similarity между узлом $B$ и query-векторами.

Тогда переходный вклад вдоль ребра будет определяться не только весом ребра, но
и тем, насколько приемник вообще подходит к смыслу запроса.

\subsection{Постановка}

Стенд сохранялся максимально близким к \texttt{exp3}:

\begin{itemize}[leftmargin=2em]
\item тот же hand-authored кофейный граф;
\item тот же запрос ``why coffee gives energy'';
\item те же embeddings;
\item единственное изменение -- наличие или отсутствие semantic gate.
\end{itemize}

Это был хороший чистый опыт: если ranking исправится, значит, проблема
действительно была только в том, что диффузия игнорировала query relevance.

\subsection{Что ожидалось при двух конкурирующих картинах}

\begin{enumerate}[leftmargin=2em]
\item \textbf{Картина A: причина в отсутствии relevance-gate.}
Тогда после добавления gate \texttt{caffeine} должен обогнать
\texttt{beverage}.
\item \textbf{Картина B: причина не только в relevance-gate, но и в самой
топологии сходимости путей.} Тогда gate улучшит картину, но не перевернет ее.
\end{enumerate}

\subsection{Результат}

Вторая картина оказалась верной. После введения gate распределение изменилось в
нужную сторону, но не настолько, чтобы исправить ranking:

\begin{center}
\begin{tabular}{lcc}
\toprule
Узел & Без gate & С gate \\
\midrule
\texttt{beverage} & 146.67 & 129.46 \\
\texttt{caffeine} & 26.67 & 27.20 \\
\texttt{juice} & 23.47 & 11.49 \\
\texttt{soda} & 23.47 & 13.65 \\
\bottomrule
\end{tabular}
\end{center}

То есть gate работал \emph{качественно правильно}, но \emph{количественно
недостаточно}.

\subsection{Почему это особенно важный результат}

Здесь произошла одна из самых тонких и самых полезных развязок всей ранней
серии. Можно было бы поспешно сказать:

\begin{quote}
``Ну вот, similarity плохая. Посмотрим на слово \texttt{energy} и увидим, что
\texttt{beverage} к нему даже ближе, чем \texttt{caffeine}. Значит, виновато
embedding-пространство''.
\end{quote}

Но аккуратный разбор показал, что это ложный след. По фактическому правилу
агрегации $\mathrm{rel}(B\mid q)$ узел \texttt{caffeine} \emph{и так} был
релевантнее \texttt{beverage}. Проигрыш происходил не на уровне per-node
relevance, а на уровне \emph{суммарной амплитуды после сходимости нескольких
путей}.

\begin{takeaway}
\texttt{exp4} показал очень важную вещь: можно иметь узел, который \emph{сам по
себе} более релевантен запросу, и все равно проигрывать, если конкурирующий узел
получает поток по нескольким сходящимся маршрутам. Это первый явный выход серии
к проблеме convergence topology.
\end{takeaway}

\subsection{Какой следующий вопрос возник немедленно}

После \texttt{exp4} стало недостаточно обсуждать только similarity. Требовался
эксперимент, который бы изолировал следующее утверждение:

\begin{quote}
``Если менять только топологию сходимости, оставляя запрос, embeddings и gate
неизменными, то ranking действительно должен переворачиваться''.
\end{quote}

Это и становится задачей \texttt{exp6}; но до него в серии появляется еще одна
ветвь -- \texttt{exp5}, посвященная не retrieval, а write-path.

\section{\texttt{exp5}: write-path как отдельный архитектурный срез}

\subsection{Почему эта глава здесь вообще стоит}

На первый взгляд \texttt{exp5} выбивается из серии. Вокруг него нет ни
\texttt{caffeine}, ни \texttt{PageRank}, ни topology correction. Но именно это и
важно: память не может состоять только из retrieval. Нужен механизм записи новых
узлов, дедупликации, коммита и форензического поиска по времени.

\texttt{exp5} поэтому играет роль боковой, но принципиально важной ветви: он
проверяет, что граф не только можно читать, но и можно \emph{корректно
модифицировать}.

\subsection{Постановка}

Сценарий был полностью контролируемым. Вместо автоматического извлечения
кандидатов из естественного текста вручную задавались объекты \texttt{Candidate},
имитирующие результат будущего экстрактора.

Сюжет -- ``поездка на Азоры'':

\begin{enumerate}[leftmargin=2em]
\item покупка плавок для поездки;
\item повтор того же события, который должен быть опознан как дубликат;
\item покупка зонта для той же поездки;
\item отдельное событие про ноутбук в другой день.
\end{enumerate}

Проверялись:

\begin{itemize}[leftmargin=2em]
\item корректный коммит узлов;
\item пропуск эпистемического дубликата;
\item forensic index по timestamp;
\item ситуация, где один кандидат ссылается на другой кандидат из того же batch.
\end{itemize}

\subsection{Какая альтернатива тут была на самом деле}

Хотя сценарий выглядит техническим, в нем скрыта содержательная архитектурная
развилка.

\begin{enumerate}[leftmargin=2em]
\item \textbf{Наивный single-pass commit.} Кандидаты коммитятся и сразу
разрешают свои ребра.
\item \textbf{Batch-aware two-pass commit.} Сначала весь пакет кандидатов
проходит gate и получает id, а уже потом разрешаются их внутренние ссылки.
\end{enumerate}

Если второй режим не нужен, значит внутренняя структура event-batch не играет
роли. Если нужен, значит write-path обязан уважать пакетность события.

\subsection{Что обнаружилось}

На первом же прогоне всплыл настоящий баг: если первый кандидат ссылается на
сущность, которую второй кандидат того же batch должен определить как полноценный
узел, single-pass логика создает stub и тем самым портит последующий commit.

В конкретном сценарии это означало, что \texttt{azores trip} оказывался не
\texttt{episode}-узлом, а generic \texttt{claim}-stub. Более того, правильный
кандидат затем ошибочно считался дубликатом уже существующего узла.

\subsection{Исправление}

Исправление состояло в переходе к двухпроходной логике:

\begin{enumerate}[leftmargin=2em]
\item весь batch проходит gating;
\item все одобренные кандидаты создают голые узлы;
\item внутри batch строится \texttt{content -> id} map;
\item только после этого разрешаются ребра, причем сначала по этой карте, и лишь
затем через механизм stub.
\end{enumerate}

\subsection{Результат после исправления}

После введения двухпроходного коммита сценарий прошел чисто:

\begin{itemize}[leftmargin=2em]
\item оба нужных узла коммитились корректно;
\item дубликат пропускался;
\item \texttt{azores trip} получал правильный \texttt{node\_kind};
\item forensic lookup по дню 9 находил только релевантные узлы и не подмешивал
ноутбук из другого дня.
\end{itemize}

\subsection{Почему этот опыт важен для всей книги}

\texttt{exp5} показывает, что у архитектуры памяти есть как минимум две разные
линии зрелости:

\begin{itemize}[leftmargin=2em]
\item retrieval-линия, где исследуется прохождение и ранжирование сигнала;
\item write-path-линия, где исследуется корректность коммита и структуры
обновления графа.
\end{itemize}

Именно благодаря этому \texttt{exp5} важно не потерять в общем своде. Без него
может возникнуть ложное впечатление, что серия занималась только чтением памяти.
На деле уже здесь проясняется, что даже простой commit требует не меньшей
аккуратности, чем retrieval.

\begin{openbox}
\texttt{exp5} оставил открытым следующий вопрос: что будет, когда ссылки начнут
идти не только внутри одного batch, но и между несколькими событиями, особенно
если ранний event сначала породил stub, а поздний должен его доопределить. Это
уже не проблема локального коммита, а проблема долговременной консолидации и
ревизии графа.
\end{openbox}

\section{Промежуточный итог после \texttt{exp4} и \texttt{exp5}}

К этому моменту исследовательская линия уже расщепилась на две зрелые оси.

\begin{enumerate}[leftmargin=2em]
\item Retrieval-ось показала, что простой semantic gate не решает проблему
слишком ярких широких узлов.
\item Write-path-ось показала, что даже controlled commit требует
batch-aware-процедуры и не выдерживает наивного single-pass режима.
\end{enumerate}

Дальнейшая серия \texttt{6x} возвращается именно к retrieval-оси и начинает
разбирать ее уже более математически: что именно удерживает сигнал, и почему
некоторые узлы светятся слишком ярко не по смыслу, а по топологии и глобальной
славе.
